% ============================================
% Section: System Model and Wireless Channel
% ============================================
\label{sec:system}

We formalize the node model, per-node multi-sensor fusion, the wireless
channel model, and the traffic characterization that drives protocol design.

\subsection{Node Model}

A node $i \in \{1,\dots,N\}$ is a mobile platform carrying:
\begin{itemize}
  \item \textbf{$N_c^{(i)} \geq 1$ cameras} with partially overlapping FOVs,
    known intrinsics $K^{(i,c)}$, distortion $D^{(i,c)}$, and resolutions
    $W^{(i,c)}\times H^{(i,c)}$;
  \item \textbf{$N_l^{(i)} \geq 1$ LiDARs} of possibly heterogeneous types
    (mechanical spinning, solid-state, MEMS), each with known beam count,
    FOV, and scan rate;
  \item \textbf{1 virtual IMU} fusing all physical IMUs into a single body-frame
    stream at 200--400~Hz;
  \item \textbf{edge GPU} (tier: Orin, 4090, A100) providing $\geq$20 TFLOPS
    for local SLAM and compression;
  \item \textbf{wireless NIC} (Wi-Fi 6/6E or 5G NR modem).
\end{itemize}

All intrinsics, extrinsics $T_{B_i\leftarrow S}$, and time offsets
$\Delta t_{S}$ relative to the IMU are calibrated offline or estimated online
(Section~\ref{sec:ctfusion}) and persisted in the signed node descriptor.

\subsection{Wireless Channel Model}

\label{sec:channel}

The wireless environment is modeled to capture the constraints that drive
protocol design. We consider two deployment scenarios:

\paragraph{Scenario A --- Wi-Fi 6/6E infrastructure.} Nodes connect to one or
more Wi-Fi~6 APs (802.11ax, 2$\times$2 MIMO, 80~MHz channels). The per-link
PHY rate $R_{\text{phy}}^{(i)}$ depends on MCS index $m^{(i)}(t)$, spatial
streams $S_{\text{ss}}$, and channel bandwidth $W$:
\begin{equation}
  R_{\text{phy}}^{(i)}(t) = S_{\text{ss}} \cdot W \cdot \eta_{\text{code}}(m^{(i)}(t)) \cdot \log_2(M(m^{(i)}(t))),
\end{equation}
where $\eta_{\text{code}}$ is the LDPC code rate and $M$ the modulation order.
The effective MAC throughput, accounting for EDCA contention, aggregation
(A-MPDU), and overhead, is:
\begin{equation}
  R_{\text{eff}}^{(i)}(t) = \alpha_{\text{MAC}} \cdot R_{\text{phy}}^{(i)}(t),
\end{equation}
with $\alpha_{\text{MAC}} \in [0.6, 0.8]$ under moderate contention. At MCS~11
(1024-QAM, 5/6 code rate), $R_{\text{eff}} \approx 480$~Mbps per link; at
MCS~0 (BPSK, 1/2) under deep fade, $R_{\text{eff}} \approx 3$~Mbps---a 160$\times$
dynamic range that the protocol must handle.

\paragraph{Scenario B --- 5G NR.} Nodes connect to gNB via NR-U (unlicensed
spectrum) or licensed NR with configured grant URLLC for pose streams and
dynamic grant eMBB for submap uploads. Per-flow QoS is enforced via 5QI values;
we map our four traffic classes to 5QI~3 (GBR, 50~ms), 5QI~7 (non-GBR, 100~ms),
5QI~8 (non-GBR, 300~ms), and 5QI~9 (non-GBR, default).

\paragraph{Link quality metric.} Each node computes a moving exponential average
of its effective throughput and packet delivery ratio (PDR):
\begin{align}
  \bar{R}_{\text{eff}}^{(i)}[k] &= \beta \cdot \bar{R}_{\text{eff}}^{(i)}[k-1] + (1-\beta) \cdot R_{\text{eff}}^{(i)}[k], \quad \beta=0.9, \\
  \text{PDR}^{(i)}[k] &= \frac{\text{pkts\_acked}[k]}{\text{pkts\_sent}[k]}.
\end{align}
The composite link quality $Q^{(i)} = \bar{R}_{\text{eff}}^{(i)} \cdot \text{PDR}^{(i)}$ feeds into
cluster-head election (Section~\ref{sec:election}) and compression budget
allocation (Section~\ref{sec:compression}).

\subsection{Design Goals}

\label{sec:design}

Co-LIC2 targets seven quantitative design goals, summarized in
Table~\ref{tab:goals_list}. These goals drive every protocol and algorithm
decision in subsequent sections and serve as the evaluation criteria in
Section~\ref{sec:experiments}.

\begin{table}[t]
  \centering
  \caption{Co-LIC2 design goals.}
  \label{tab:goals_list}
  \footnotesize
  \setlength{\tabcolsep}{3pt}
  \begin{tabular}{@{}p{0.4cm}p{1.6cm}p{5.0cm}@{}}
    \toprule
    \# & Goal & Target \\
    \midrule
    G1 & Scalability & 2--16 nodes; sub-linear per-link BW growth \\
    G2 & Link-loss robustness & map quality within 5\% of connected after 30 s outage \\
    G3 & Global consistency & ATE $<$0.3 m on 200 m loops with 3 nodes \\
    G4 & Photo-realism & PSNR within 1 dB of single-node GLIC2 baseline \\
    G5 & Real-time & local odometry $\ge$20 Hz; submap merge $\le$2 s latency \\
    G6 & Heterogeneity & nodes may differ in sensor count/quality \\
    G7 & Backward compat. & a single node with coop off = Gaussian-LIC2 \\
    \bottomrule
  \end{tabular}
\end{table}

\subsection{Per-Node Continuous-Time Multi-Sensor Fusion}

\label{sec:ctfusion}

Each node estimates a non-uniform cubic B-spline trajectory $\Twb^{(i)}(t) \in \SEthree$
using a generalized version of Coco-LIC~\cite{lang2023coco} and Gaussian-LIC2~\cite{lang2026gaussian2}.
Control points are placed adaptively based on IMU acceleration magnitude.

\paragraph{IMU residuals.} For measurement at $t_k$:
\begin{equation}
  \mathbf{r}_{\text{IMU}}(t_k) =
  \begin{bmatrix} \mathbf{a}_k - \hat{\mathbf{a}}(t_k) \\ \boldsymbol{\omega}_k - \hat{\boldsymbol{\omega}}(t_k) \end{bmatrix}_{\Sigma_{\text{IMU}}^{-1}},
\end{equation}
where $\hat{\mathbf{a}}$, $\hat{\boldsymbol{\omega}}$ derive from the spline's
analytical derivatives.

\paragraph{LiDAR residuals.} For each LiDAR $l$, point $\mathbf{p}_k^{(l)}$ at $t_k$:
\begin{equation}
  \mathbf{r}_{\text{LiDAR}}^{(l)}(t_k) = \mathbf{n}^\top \bigl( \Twb^{(i)}(t_k)\, T_{B_i\leftarrow L_l}\, \mathbf{p}_k^{(l)} - \boldsymbol{\mu}_* \bigr),
\end{equation}
with $\boldsymbol{\mu}_*$ the nearest Gaussian mean in $\GaussianMap_i$, $\mathbf{n}$ the
local surface normal from Gaussian covariances.

\paragraph{Camera residuals.} For each camera $c$, keyframe at $t_k$: SSIM+L1
photometric residual against the rendered $\GaussianMap_i$ at spline-evaluated
pose, exactly as in Gaussian-LIC2.

\paragraph{Time-offset residuals.} $N_c + N_l$ unknown offsets relative to IMU
are estimated jointly with the spline in a sliding-window Gauss-Newton solve.
The full residual stack with B-spline motion prior yields the per-node
trajectory and local Gaussian map $\GaussianMap_i$.

\subsection{Traffic Characterization}

\label{sec:traffic}

Co-LIC2 generates four traffic classes with distinct QoS requirements:

\begin{table}[t]
  \centering
  \caption{Traffic classes and QoS requirements.}
  \label{tab:traffic}
  \footnotesize
  \setlength{\tabcolsep}{3pt}
  \begin{tabular}{@{}p{1.3cm}p{1.6cm}p{1.0cm}p{1.3cm}p{0.9cm}p{0.8cm}p{1.0cm}@{}}
    \toprule
    Class & Dir.\@ & Rate & Pkt size & Latency & PDR & Streams \\
    \midrule
    \textbf{Pose}     & node$\to$head & 20 Hz  & 200 B    & $<$50 ms  & $>$99\%   & 1/node \\
    \textbf{Submap}   & node$\to$head & 0.33 Hz& 0.5--5 MB& $<$2 s    & $>$95\%   & 1/node \\
    \textbf{Keyframe} & node$\to$head & 1 Hz   & 200 KB   & $<$500 ms & $>$99\%   & 1--4/node \\
    \textbf{Control}  & bidir.\@       & 1--5 Hz& 500 B    & $<$100 ms & $>$99.9\% & 1/node \\
    \bottomrule
  \end{tabular}
\end{table}

\textbf{Pose stream} carries B-spline control-point updates for real-time
visualization and early registration. \textbf{Submap} is the bulk payload:
a compressed spatial sub-volume of Gaussians with user-defined temporal/spatial
boundaries. \textbf{Keyframe} carries compressed reference images and NetVLAD
descriptors for visual loop closure. \textbf{Control} handles HELLO, HEARTBEAT,
ELECTION, CONSTRAINT, MAP\_DELTA, LEAVE.

This multi-modal, heterogeneous traffic mix is the motivation for the
QoS-aware, multi-stream protocol design of Section~\ref{sec:protocol}.
